\begin{abstract}
Prime-editing efficiency predictors are commonly compared by correlation across measured designs, whereas their practical use requires choosing among alternative pegRNAs for a fixed intended edit. We investigated the relationship between these tasks using \model{}, a sequence model combining paired wild-type/edited representations, context conditioning and ordinal or proportion-based supervision. On the 20,509-row source benchmark, its frozen heterogeneous ensemble achieved Spearman correlation 0.9079 versus 0.8690 for released OptiPrime. Corrected fixed-edit evaluation also favored \model{} on 1,463 informative source decisions: the mean efficiency of the highest-scored candidate was 0.13636 versus 0.12721. However, on an initially reserved external library containing 30,475 eligible decisions, the source-trained ordinal--state-space ensemble retained higher correlation (0.7091 versus 0.6931) but lower top-choice efficiency (0.03627 versus 0.03690; paired difference $-0.00064$, 95\% interval $[-0.00089,-0.00038]$). Audits identified consequential distinctions between decision grouping and batch composition, prediction-head retention and deployed-selector retention, and optimization-seed and label-acquisition replication. Target-supervised ordinal fine-tuning improved top-choice efficiency to 0.04178 versus 0.04020 for unadapted OptiPrime on a separate retrospective library test subset. In 76 subsequent computational fits, stronger optimization controls and repeated label acquisitions removed a simple source-anchoring advantage. A source-balanced teacher-margin constraint improved retention, but constrained adaptation sacrificed target performance and did not outperform OptiPrime at a nominal 1,000-group label budget. These results establish a dataset-specific prediction improvement and a supervised adaptation benefit, while showing why neither implies universally superior design selection. The resulting computational framework evaluates actual deployed decisions, label cost, source retention and evaluation exposure together; independent-study confirmation remains outstanding.
\end{abstract}
